\documentclass[aspectratio=169]{beamer}
\usetheme{Madrid}
\usecolortheme{seagull}
\usepackage[T1]{fontenc}
\usepackage{lmodern}
\usepackage{amsmath}
\usepackage{booktabs}

\title{FerrumMD: Molecular Dynamics in Rust}
\subtitle{Main components for Rust developers new to MD}
\author{FerrumMD project overview}
\date{\today}

\begin{document}

\begin{frame}
  \titlepage
\end{frame}

\begin{frame}{Why this project matters}
\begin{itemize}
  \item \textbf{Goal:} build an MD engine with Rust's safety + performance model.
  \item \textbf{Audience bridge:} map familiar systems ideas (data layout, kernels, parallelism) to simulation workflows.
  \item \textbf{Scope today:} short-range interactions, molecular topologies, integrators, thermodynamic controls, and I/O.
\end{itemize}
\end{frame}

\begin{frame}{MD in one minute}
\begin{itemize}
  \item MD evolves particle positions/velocities by integrating Newton's equations.
  \item Core loop per step:
  \begin{enumerate}
    \item compute forces from current geometry,
    \item update positions/velocities,
    \item apply boundary conditions and controls (thermostat/barostat),
    \item report observables and save trajectory.
  \end{enumerate}
  \item Think of it as a repeated \texttt{state -> forces -> state'} transform with strict numerical constraints.
\end{itemize}
\end{frame}

\begin{frame}{High-level FerrumMD architecture}
\begin{block}{Data and state}
\texttt{state}, \texttt{molecule}, \texttt{parameters}, \texttt{cell}, \texttt{cell\_subdivision}
\end{block}
\begin{block}{Physics kernels}
\texttt{lennard\_jones}, bonded interactions, electrostatics roadmap, free-energy and quantum modules
\end{block}
\begin{block}{Control and workflows}
\texttt{ensembles} (NVE/NVT/NPT), \texttt{thermostat\_barostat}, minimization, replica exchange
\end{block}
\begin{block}{Interfaces}
CLI binaries, GRO/XTC trajectory output, optional viewer feature, optional Python + MPI features
\end{block}
\end{frame}

\begin{frame}{Core simulation state (Rust view)}
\begin{itemize}
  \item Struct-based system representation for particles, molecules, box, and parameters.
  \item Explicit unit convention (nm, ps, amu, kJ/mol) avoids hidden conversion bugs.
  \item Rust typing helps separate concerns:
  \begin{itemize}
    \item immutable parameter sets,
    \item mutable evolving state,
    \item module boundaries for force kernels and integration.
  \end{itemize}
  \item Result: easier reasoning about correctness and reproducibility.
\end{itemize}
\end{frame}

\begin{frame}{Force models implemented}
\begin{columns}[T]
\column{0.52\textwidth}
\textbf{Nonbonded}
\begin{itemize}
  \item Lennard--Jones pair interaction
  \item Minimum-image convention with periodic boundary conditions
  \item Cell subdivision module for scaling pair computations
\end{itemize}

\column{0.45\textwidth}
\textbf{Bonded}
\begin{itemize}
  \item Harmonic bond forces
  \item Support for small molecular systems (e.g., H$_2$)
  \item Topology readers for Martini/CHARMM-style inputs
\end{itemize}
\end{columns}
\end{frame}

\begin{frame}{Time integration and thermodynamic control}
\begin{itemize}
  \item Velocity-Verlet integrator as the default engine.
  \item Ensemble support:
  \begin{itemize}
    \item NVE (energy-conserving baseline),
    \item pseudo-NVT/NVT with thermostats,
    \item NPT path with isotropic barostat support.
  \end{itemize}
  \item Thermostats implemented: Berendsen and Nose--Hoover.
  \item Initialization: Maxwell--Boltzmann velocity sampling.
\end{itemize}
\end{frame}

\begin{frame}{I/O and ecosystem interoperability}
\begin{itemize}
  \item Reads molecular inputs: PDB, GRO, LAMMPS data/dump.
  \item Force-field conversion helpers: Martini and CHARMM parsers.
  \item Writes GRO/XTC outputs for external visualization (e.g., VMD).
  \item Includes a Rust-native viewer scaffold (feature-gated) for extension in pure Rust.
  \item Python bindings scaffold and MPI-enabled path for integration and parallel demos.
\end{itemize}
\end{frame}

\begin{frame}{Why Rust is interesting here}
\begin{itemize}
  \item \textbf{Safety:} ownership/borrowing reduces hidden aliasing in mutable simulation state.
  \item \textbf{Performance:} predictable data-oriented code, low abstraction penalty.
  \item \textbf{Composability:} feature flags (MPI/viewer/python) keep core minimal.
  \item \textbf{Reliability culture:} strong testing and explicit errors match scientific reproducibility goals.
\end{itemize}
\end{frame}

\begin{frame}{Roadmap and open opportunities}
\begin{itemize}
  \item Rigid-water constraints (SETTLE/SHAKE-RATTLE path).
  \item Topology-correct nonbonded exclusions/pairs.
  \item Electrostatics and units audit for production TIP3P fidelity.
  \item Possible Rust-focused contributions:
  \begin{itemize}
    \item SIMD kernel specialization,
    \item domain decomposition beyond demo MPI,
    \item property-based tests for invariants,
    \item GPU offload experimentation.
  \end{itemize}
\end{itemize}
\end{frame}

\begin{frame}{Takeaways for Rust developers}
\begin{enumerate}
  \item FerrumMD already demonstrates an end-to-end MD stack in Rust.
  \item The project is a strong playground for systems + scientific computing ideas.
  \item You can contribute without deep prior MD experience by starting from architecture and tooling modules.
\end{enumerate}
\vspace{0.8em}
\centering
\large Questions \& discussion
\end{frame}

\end{document}
